\begin{textsample}{NEG Dim 1 – 2015\_c4 – Score -4.00 – 2015\_c4/enade\_2015\_teo\_839aebf29e.txt}  \label{ex:f1_neg_103}
CAPÍTULO 4 PERCEPÇÃO DA PROVA As análises feitas neste capítulo tratam das percepções dos concluintes da Área de Teologia sobre a prova aplicada no Enade/2015.
Essas percepções foram mensuradas por meio de nove questões que avaliaram desde o grau de dificuldade da prova até o tempo gasto para concluí -la.
As percepções sobre a prova foram relacionadas com o desempenho dos estudantes e com a Grande Região de funcionamento do curso.
O questionário de percepção da prova encontra -se ao final do Anexo VII, que traz a reprodução da prova.
O desempenho dos estudantes foi classificado em quatro quartos.
Para tanto, esse desempenho foi ordenado de forma ascendente.
O percentil 25, P25, também conhecido como primeiro quartil, é a nota de desempenho que deixa um quarto ( 25 \% ) dos valores observados abaixo e três quartos acima.
A Figura 1 apresenta uma ilustração deste conceito.
O quarto inferior de desempenho é composto pelas notas abaixo do primeiro quartil.
Já o percentil 75, P75, também conhecido como terceiro quartil, é o valor para o qual há três quartos ( 75 \% ) dos dados abaixo e um quarto acima dele.
O quarto superior de desempenho é composto pelas notas iguais ou acima do terceiro quartil.
O percentil 50, P50, também conhecido como mediana, é o valor que divide as notas em dois conjuntos de igual tamanho.
O segundo quarto inclui valores entre o primeiro quartil ( P25 ) e a mediana.
O terceiro quarto contém os valores entre a mediana ( P50 ) e o terceiro quartil ( P75 ).
Vale ressaltar que percentis, quartis e medianas são pontos que não obrigatoriamente pertencem ao conjunto original de dados, ao passo que os quartos são subconjuntos dos dados originais.
A seguir, serão apresentados gráficos com resultados selecionados, relativos às nove questões avaliadas por grupos de estudantes.
Os gráficos apresentam nas barras o percentual de alunos que assinalaram uma das opções ou a soma das porcentagens daqueles que assinalaram duas ( ou três ) delas.
Por exemplo, para as questões 1 e 2, os gráficos apresentam a porcentagem total de participantes que assinalaram as opções ( D ) Difícil e ( E ) Muito difícil.
Em \textbf{cada} barra foram assinalados também os extremos do intervalo de confiança de 95 \% como linhas verticais unidas por uma linha horizontal na forma da letra H maiúscula.
As Tabelas no Anexo II apresentam os valores absolutos e a distribuição percentual17 das alternativas válidas das nove questões, segundo o mesmo recorte de desempenho dos alunos e Grande Região de funcionamento do curso. 4.1 GRAU DE DIFICULDADE DA PROVA 4.1.1 Componente de Formação Geral Ao avaliarem “ Qual o grau de dificuldade desta prova na parte de Formação Geral? ” ( Questão 1 ), 34,2 \% do grupo de inscritos e presentes optaram pelas alternativas Difícil ou Muito difícil.
Entretanto, para mais da metade dos estudantes ( 56,4 \% ), o Componente de Formação Geral da prova foi considerado com grau de dificuldade Médio ( Gráfico 4.1, Gráfico 4.2 e, no Anexo II, a Tabela II.1 ).
O percentual de estudantes que consideraram a prova como Difícil ou Muito difícil foi maior na região Centro-Oeste, onde a proporção foi de 41,1 \%, enquanto a de menor incidência foi a Sudeste, com 31,6 \%.
No Gráfico 4.1, é possível observar que há diferença estatisticamente significativa entre as regiões Sudeste e Sul ( 37,8 \% ).
Nas Grandes Regiões, a proporção de presentes à prova que consideraram o Componente de Formação Geral como sendo de grau de dificuldade Médio \textbf{esteve} entre 54,0 \% na região Nordeste e 57,6 \% na região Sudeste. 17 Cumpre \textbf{lembrar} uma das convenções para tabelas numéricas ( pág. iii ) sobre a possibilidade de a soma das partes não resultar em 100 \% por questões de arredondamento.
O percentual de alunos que consideraram a prova Difícil ou Muito difícil foi decrescente em função dos quartos de desempenho : 42,6 \% no primeiro quarto e 22,2 \% no quarto, grupo de melhor desempenho na prova.
Nos quartos de desempenho intermediários, a proporção de alunos que consideraram a prova Difícil ou Muito difícil foi de 38,5 \% no 2º quarto e 34,6 \% no 3º quarto.
Há diferenças estatisticamente significativas entre o quarto de desempenho superior e todos os demais, assim como entre o primeiro e o terceiro.
Para todos os quartos de desempenho, a alternativa modal para esta pergunta foi Médio, com 48,4 \% e 64,4 \% dos respondentes nos quartos extremos, primeiro e quarto, respectivamente. 4.1.2 Componente de Conhecimento Específico Ao responderem à Questão 2 – “ Qual o grau de dificuldade desta prova na parte de Componente Específico? ” – 37,7 \% do grupo de estudantes classificaram -na como Difícil ou Muito difícil.
O Componente de Conhecimento Específico da prova foi considerado com grau de dificuldade Médio por 53,1 \% dos alunos ( Gráfico 4.3, Gráfico 4.4, e, no Anexo II, a Tabela II.2 ).
A análise das respostas dos estudantes quanto ao grau de dificuldade do Componente de Conhecimento Específico da prova, agregado por Grande Região, mostra que a diferença entre a maior e a menor proporção de alunos que a avaliaram como Difícil ou Muito difícil é estatisticamente significativa : a maior na região Norte ( 46,0 \% ) e a menor proporção na região Sudeste ( 36,0 \% ).
O percentual de alunos que classificaram o grau de dificuldade como Médio, no Componente de Conhecimento Específico, variou de 48,8 \% a 56,0 \%, para as regiões Norte e Centro-Oeste, respectivamente.
Considerando -se a avaliação da dificuldade das questões do Componente de Conhecimento Específico da prova, de acordo com o desempenho dos estudantes, observa -se que há diferença estatisticamente significativa do resultado do quarto de melhor desempenho em relação aos demais.
A proporção dos que classificaram a parte específica como Difícil ou Muito difícil variou de 28,6 \% ( último quarto ) a 43,1 \% ( primeiro quarto ).
A alternativa modal para a Questão 2 foi o grau Médio, com 46,3 \% do quarto inferior e 60,4 \% do superior optando por essa resposta. 4.2 EXTENSÃO DA PROVA EM RELAÇÃO AO TEMPO TOTAL Indagados quanto à extensão da prova, em relação ao tempo total oferecido para a sua resolução ( Questão 3 ), os estudantes apontaram, com maior incidência, a alternativa que considerava a extensão adequada, para todas as agregações consideradas ( Gráfico 4.5, Gráfico 4.6, e, no Anexo II, a Tabela II.3 ).
O percentual de alunos que responderam ser a extensão da prova adequada foi de 41,1 \%.
Já 51,0 \% dos inscritos presentes consideraram que a prova foi longa ou muito longa e 7,8 \% a avaliaram como curta ou muito curta.
Entre as Grandes Regiões a proporção daqueles que avaliaram a prova como longa ou muito longa em relação ao tempo total destinado à sua resolução variou de 42,9 \% na região Norte até 54,6 \% na região Centro-Oeste.
A diferença entre as regiões Norte e Sudeste é estatisticamente significativa.
Considerando -se o desempenho dos alunos, nota -se ainda que 38,9 \% consideraram a extensão da prova adequada no quarto de desempenho inferior e 46,2 \% no de melhor desempenho ( quarto superior ).
A proporção de estudantes que consideraram a prova longa ou muito longa em função dos quartos de desempenho variou de 49,0 \% ( quarto superior ) a 52,1 \% ( quartos intermediários ).
No quarto inferior a proporção que escolheu essa categoria foi 51,0 \%.
No Gráfico 4.6, constata -se que não há diferença estatisticamente significativa entre as proporções de escolha desta resposta. 4.3 COMPREENSÃO DOS ENUNCIADOS DAS QUESTÕES 4.3.1 Componente de Formação Geral Com relação aos enunciados das questões do Componente de Formação Geral ( Questão 4 ), as opiniões foram positivas, já que 68,5 \% dos alunos avaliados consideraram todos ou a maioria dos enunciados das questões “ claros e objetivos ” ( Gráfico 4.7, Gráfico 4.8, e, no Anexo II, a Tabela II.4 ).
Na análise regional, a percentagem de estudantes que avaliaram que todos ou a maioria dos enunciados das questões do Componente de Formação Geral \textbf{estavam} “ claros e objetivos ” variou de 59,6 \% na região Norte a 71,5 \% na região Sudeste, sendo essa diferença e a diferença entre Nordeste e Sudeste, as únicas diferenças estatisticamente significativas.
A análise das percepções dos estudantes sobre a clareza e objetividade dos enunciados permite afirmar que todos, ou a maioria dos enunciados de questões relativas ao Componente de Formação Geral, foram considerados “ claros e objetivos ” para a maior parte dos respondentes ( maior do que 62 \% em todas as regiões e maior do que 56 \% para todos os quartos de desempenho ).
Segundo o desempenho, observa -se que a proporção dos que emitiram esta opinião cresce conforme o desempenho aumenta, com diferenças estatisticamente significativas do quarto de melhor desempenho ( 78,5 \% ) para os demais e do primeiro ( 56,6 \% ) para o segundo ( 67,0 \% ) e terceiro ( 70,6 \% ) quartos. 4.3.2 Componente de Conhecimento Específico Com relação aos enunciados das questões do Componente de Conhecimento Específico da prova, para 68,8 \% dos estudantes avaliados da Área de Teologia, a clareza e a objetividade ( Questão 5 ) \textbf{estavam} presentes em todas ou na maioria das questões ( Gráfico 4.9, Gráfico 4.10, e no Anexo II, a Tabela II.5 ).
A maioria dos estudantes de todas as Grandes Regiões brasileiras considerou “ claros e objetivos ” todos ou a maioria dos enunciados das questões do Componente de Conhecimento Específico da prova, percentual sempre maior ou igual a 62,6 \% ( regiões Norte e Nordeste ).
A proporção dessa resposta nas regiões Sudeste e Sul foram, respectivamente, 69,8 \% e 71,9 \%, havendo diferença estatisticamente significativa dessas duas regiões em relação à região Nordeste.
A proporção de estudantes que consideraram os enunciados das questões “ claros e objetivos ” apresenta uma tendência crescente em relação ao aumento de desempenho : mais elevada no quarto superior ( 78,1 \% ) se comparada ao quarto inferior de desempenho ( 55,6 \% ).
As diferenças do quarto superior para os demais quartos e do primeiro para os dois seguintes são estatisticamente significativas. 4.4 SUFICIÊNCIA DAS INFORMAÇÕES/INSTRUÇÕES FORNECIDAS Ao avaliarem as informações/instruções fornecidas para a resolução das questões ( Questão 6 ), 81,4 \% dos respondentes da Área de Teologia de todo o Brasil afirmaram que estas eram até excessivas ou “ suficientes ” em todas ou na maioria das questões ( Gráfico 4.11, Gráfico 4.12, e, no Anexo II, a Tabela II.6 ).
Quanto à distribuição de respondentes pelas Grandes Regiões observa -se que a proporção de estudantes que consideraram as informações/instruções fornecidas até excessivas ou “ suficientes ” em todas ou na maioria das questões foi sempre superior a 76 \%, chegando a 85,6 \% na região Sul.
As diferenças entre o maior percentual, obtido na região Sul, e os menores, das regiões Nordeste ( 76,1 \% ) e Norte ( 76,7 \% ) são estatisticamente significativas.
Levando -se em conta o desempenho dos participantes, notam -se diferenças estatisticamente significativas entre as opiniões de estudantes do quarto superior de desempenho em relação aos demais, o mesmo ocorrendo no quarto inferior em relação aos intermediários, como mostra o Gráfico 4.12.
O percentual de participantes que avaliaram “ as informações/instruções ” como até excessivas ou “ suficientes ” em todas ou na maioria das questões foi mais elevado no quarto superior ( 90,9 \% ), percentual superior à média nacional ( 81,4 \% ).
Já no quarto inferior, a suficiência das informações/instruções declarada como até excessiva, em todas ou na maioria das questões foi percebida por 69,3 \% dos respondentes.
Estas proporções são crescentes com os quartos de desempenho. 4.5 DIFICULDADE ENCONTRADA AO RESPONDER À PROVA \textbf{Perguntados} sobre as dificuldades com as quais se depararam ao responder à prova ( Questão 7 ), 8,4 \% dos estudantes apontaram o Desconhecimento do conteúdo.
Para 56,5 \%, a Forma diferente de abordagem do conteúdo foi indicada como dificuldade.
Já a Falta de motivação para fazer a prova foi a dificuldade apontada por 9,0 \% dos respondentes.
Considerando -se todo o Brasil, 21,0 \% dos respondentes afirmaram que não tiveram qualquer tipo de dificuldade para responder à prova ( Tabela II.7 no Anexo II ).
Os Gráficos 4.13 e 4.14 apresentam os percentuais de estudantes que apontaram o Desconhecimento do conteúdo como dificuldade percebida ao responder à prova.
Na análise por Grandes Regiões, o percentual de inscritos e presentes que apontaram o Desconhecimento do conteúdo como dificuldade ao responder à prova não superou 11,6 \%.
Os percentuais variaram de 7,2 \% na região Norte a 11,6 \% na Centro-Oeste.
A Forma diferente de abordagem do conteúdo foi a escolha modal dos estudantes, com percentuais que variaram de 55,4 \% ( região Sudeste ) a 60,1 \% ( Centro-Oeste ).
O percentual de alunos que citaram a Falta de motivação como dificuldade variou de 7,2 \% ( região Centro-Oeste ) a 10,6 \% ( região Norte ).
Os que declararam não ter qualquer tipo de dificuldade para responder à prova variaram de 16,7 \% na região Centro-Oeste a 23,0 \% na Sudeste.
Com relação aos quartos de desempenho, o Desconhecimento do conteúdo foi a opção escolhida por 4,7 \% dos estudantes do quarto superior e 12,3 \% do quarto inferior.
Observam -se diferenças estatisticamente significativas entre o quarto superior e o primeiro e o segundo ( 9,6 \% ) quartos, assim como entre o primeiro quarto e o terceiro ( 7,3 \% ).
A alternativa modal para os alunos, quando agregados pelos quartos de desempenho, foi que a dificuldade encontrada foi causada pela Forma diferente de abordagem do conteúdo : 51,5 \% no quarto inferior e 53,4 \% do quarto superior assim o responderam.
O percentual de alunos que citaram a Falta de motivação como dificuldade variou de 6,8 \% no quarto superior a 13,0 \% no quarto inferior.
Os que declararam não ter qualquer tipo de dificuldade para responder à prova variaram de 15,7 ( quarto inferior ) a 30,0 \% ( quarto superior ). 4.6 CONTEÚDOS DAS QUESTÕES OBJETIVAS DA PROVA Ao analisarem os conteúdos das questões objetivas da prova ( Questão 8 ), um percentual muito pequeno dos estudantes avaliados, apenas 8,6 \%, afirmou que não estudou ainda a maioria desses conteúdos ( Gráficos 4.15, Gráfico 4.16, e a Tabelas II.8 no Anexo II ).
A maioria ( 67,5 \% ) afirmou ter estudado e aprendido muitos ou todos os conteúdos avaliados.
Na análise por Grande Região, a proporção de respondentes que escolheram a opção não estudou ainda a maioria desses conteúdos foi pequena.
Observa -se que nas regiões Norte ( 16,2 \% ), Nordeste ( 12,7 \% ) e Centro-Oeste ( 10,2 \% ) as proporções foram maiores do que a média nacional ( 8,6 \% ).
Observa -se diferença estatisticamente significativa entre algumas das regiões.
Em todas as regiões, a maioria dos presentes afirmou ter estudado e aprendido muitos ou todos os conteúdos, com proporções variando entre 52,5 \% na região Norte e 70,5 \% na Sudeste.
Considerando -se separadamente as opiniões de estudantes dos quatro quartos de desempenho, observa -se que, no quarto inferior, 17,7 \% ofereceram como resposta que não estudou ainda a maioria desses conteúdos, sendo 2,4 \% os do quarto superior com a mesma resposta.
A diferença entre os alunos que optaram por esse motivo de dificuldade no quarto superior é estatisticamente significativa em relação aos demais quartos, assim como do quarto inferior em relação aos intermediários.
Tendo em conta o quarto superior, 87,1 \% dos alunos afirmaram ter estudado e aprendido muitos ou todos os conteúdos.
No outro extremo, no primeiro quarto, 48,8 \% optaram pelas mesmas categorias. 4.7 TEMPO GASTO PARA CONCLUIR A PROVA Ao responderem sobre o tempo de conclusão da prova ( Questão 9 ), mais da metade dos estudantes ( 75,1 \% ) afirmou ter gasto Entre duas e quatro horas ( Gráfico 4.17, Gráfico 4.18 e, no Anexo II, a Tabela II.9 ).
Considerando -se as cinco Grandes Regiões brasileiras, o percentual dos que utilizaram Entre duas e quatro horas para finalizar a prova foi superior ao percentual nacional nas regiões Sudeste ( 76,3 \% ) e Sul ( 79,8 \% ).
Nas demais Grandes Regiões o percentual de alunos que dispensaram Entre duas e quatro horas para concluir a prova, ficou entre 64,9 \% ( região Norte ) e 69,1 \% ( região Centro-Oeste ), como mostra o Gráfico 4.17.
Uma vez sendo analisadas as alternativas escolhidas pelos estudantes que se situam nos diferentes quartos de desempenho, observa -se uma tendência crescente : uma maior proporção de participantes no quarto superior ( 83,9 \% ) declarou ter gasto Entre duas e quatro horas “ para concluir a prova ” quando comparadas com os dos quartos inferiores.
As diferenças entre o primeiro ( 62,6 \% ) e os demais, bem como entre o segundo ( 73,4 \% ) e o quarto de maior desempenho são estatisticamente significativas.

% matched lemmas: cada, estar, lembrar, perguntar
\end{textsample}
